\topsection[sec:scoping]{Scoping and problem definition}

This section is the summary of phase~1. The full scoping note is \repofile{docs/01_scoping_note.md}.

\subsection{Stakes for an enterprise}
An enterprise that deploys a model it did not entirely create needs three things: the integrity of what it
deploys, the ability to show where a model came from after an incident, and a correct understanding of what a
signed model does and does not guarantee. Traditional controls (code signing, software bills of materials,
build provenance) cover source code and binaries. The behaviour of a model is set by data and by a training
process that is rarely reproducible, and those controls do not capture that lineage.

\subsection{Assets and threats}
Table~\ref{tab:assets} lists the assets the brief asks us to identify, and Figure~\ref{fig:when} places the four
threat classes of the brief (data poisoning, backdoor injection, model substitution, pipeline compromise) on the
pipeline. The distinction that organises the whole project is \emph{when} an attack acts. After an artifact is
signed, any change is in reach of cryptography. Before it is signed, a legitimate signer attests the malicious
content in good faith.

\begin{table}[H]
\caption{Assets of an AI supply chain and why each one matters.}
\label{tab:assets}
\tablefont
\begin{tabularx}{\textwidth}{P{5.4cm}L}
\toprule
\headrow\hd{Asset} & \hd{Why it matters} \\
\midrule
Datasets (training and fine-tuning) & They determine the behaviour of the model; the main target of poisoning \\
Pretrained base models & Downloaded from third parties; may carry a backdoor or executable code \\
Fine-tuned model weights & The behaviour that is deployed \\
Training and fine-tuning scripts, parameters and environment & They define how weights are produced; a compromised script or host alters the result \\
Inference and verification APIs, deployment configuration & They serve the deployed model and decide what is trusted \\
Signing keys, trust store and policy & They decide which signatures count; their theft defeats the scheme \\
Signed manifests, provenance graph and audit trail & The evidence used at deployment and after an incident \\
\bottomrule
\end{tabularx}
\end{table}

\begin{figure}[!htb]
\centering
\begin{tikzpicture}[x=1cm,y=1cm,font=\sffamily\fontsize{7.7}{9.6}\selectfont,
  art/.style={draw=cgnavy,line width=0.7pt,rounded corners=2.5pt,minimum width=2.95cm,minimum height=1.05cm,align=center,fill=white},
  atk/.style={align=center,text width=3.25cm,inner sep=2pt,text=cgink},
  flow/.style={-{Latex[length=2mm,width=1.6mm]},line width=0.8pt,color=cgnavy},
  hit/.style={-{Latex[length=1.7mm,width=1.4mm]},line width=0.6pt,color=cgmuted}]
% bands
\fill[cgband,rounded corners=3pt] (-1.75,1.02) rectangle (13.25,2.95);
\fill[cgband,rounded corners=3pt] (-1.75,-2.95) rectangle (13.25,-1.02);
\node[anchor=north west,font=\sffamily\bfseries\fontsize{8}{10}\selectfont,text=cgnavy] at (-1.6,2.9) {Before signing: the signer attests the content in good faith. Cryptography cannot see it.};
\node[anchor=south west,font=\sffamily\bfseries\fontsize{8}{10}\selectfont,text=cgnavy] at (-1.6,-2.9) {After signing: every change breaks a hash, a signature or a link. Cryptography detects it.};
% artifacts
\node[art] (a1) at (0,0) {\textbf{Datasets}\\stage 1};
\node[art] (a2) at (3.83,0) {\textbf{Base model}\\stage 2};
\node[art] (a3) at (7.66,0) {\textbf{Fine-tuned model}\\stage 3};
\node[art] (a4) at (11.49,0) {\textbf{Deployment}\\stage 4};
\draw[flow] (a1) -- (a2);
\draw[flow] (a2) -- (a3);
\draw[flow] (a3) -- (a4);
% before signing
\node[atk] (b1) at (0,1.82) {Poisoned rows in the\\training or fine-tuning data};
\node[atk] (b2) at (3.83,1.82) {Backdoored pretrained model\\from an outside publisher};
\node[atk] (b3) at (7.66,1.82) {Training host that alters\\the weights it will sign};
\node[atk] (b4) at (11.49,1.82) {Code that runs on load\\(pickle formats)};
\foreach \i in {1,2,3,4} \draw[hit] (b\i) -- (a\i);
% after signing
\node[atk] (c1) at (0,-1.82) {Rows changed, added\\or removed in storage};
\node[atk] (c2) at (3.83,-1.82) {Model file replaced\\or one weight changed};
\node[atk] (c3) at (7.66,-1.82) {Parent link rewritten;\\signature forged or replayed};
\node[atk] (c4) at (11.49,-1.82) {Revoked key or artifact used;\\trust records edited};
\foreach \i in {1,2,3,4} \draw[hit] (c\i) -- (a\i);
\end{tikzpicture}
\caption{Where attacks act along an AI pipeline. Experiments 1 to 5 of Section~\ref{sec:eval} cover the lower
band, experiments 6 and 7 the upper band. Formats that run code on load sit between the two: the file is exactly
what was signed, so the gateway refuses such formats by rule.}
\label{fig:when}
\end{figure}

\subsection{Scope, assumptions and constraints}
\textbf{In scope.} SHA-256 hashing, Ed25519 signatures, signed JSON manifests, hash-linked provenance and
dataset Merkle trees, signer roles and a stage policy, revocation, a deployment gateway, a small PyTorch
pipeline, attack scenarios and measurements, a FastAPI backend, a Streamlit dashboard, Docker and GitHub
Actions.

\textbf{Out of scope, discussed but not built.} Blockchain anchoring (optional in the brief; Section~\ref{sec:soa}
explains the choice of a signed log instead), hardware attestation, key storage in an HSM or a KMS,
reproducible training, large models, production authentication for the API, and research on new backdoor
detectors.

\textbf{Assumptions.} The signing keys, the policy file and the verifier code are trusted. Everything else
(files, storage, the index database, logs, training hosts, outside publishers, the network) may be attacked.
Each role has its own key. The held-out test split is trusted. A small model on the scikit-learn digits dataset
is enough to measure integrity guarantees and the effect of a backdoor, because the mechanisms do not depend on
the size of the model.

\textbf{Constraints.} Eight weeks, four students and no budget: open-source tools only and training on CPU.
Every result must be reproducible from the repository. What cryptography proves and what it cannot prove are
kept apart, in the code and in the writing.

\subsection{Objectives and success criteria}
The scoping note fixed five measurable objectives in the first two weeks. Table~\ref{tab:criteria} repeats
them with the result measured at the end of the project.

{\tablefont
\begin{xltabular}{\textwidth}{P{3.1cm}P{4.5cm}L}
\caption{Objectives of the scoping note, their success criteria and the measured result.}\label{tab:criteria}\\
\toprule
\headrow\hd{Objective} & \hd{Success criterion} & \hd{Result} \\
\midrule
\endfirsthead
\toprule
\headrow\hd{Objective} & \hd{Success criterion} & \hd{Result} \\
\midrule
\endhead
\bottomrule
\endfoot
Attest artifacts across at least three stages & Datasets, base model, fine-tuned model and deployment record signed and hash-linked & Met: four stages attested and linked \\
Detect tampering after attestation & Every tampering, substitution, provenance and signature variant \st{FAILED}; controls \st{VERIFIED} & Met: 48 of 48 variants ended as predicted; every manipulation after signing was rejected \\
Show the limit of attestation & Signed poisoned data verifies; backdoor success rate reported & Met: \st{VERIFIED} and deployed with 77.2\,\% (E2), 51.2\,\% (E4) and 100\,\% (E7-F) \\
Keep the cost negligible & Attestation under 5\,\% of training time; verification under 100\,ms for the chain & Met: 0.83\,\% of training time; 8.5\,ms per verification (95th percentile 12.2\,ms) \\
Reproducible system & \code{docker compose up} and the CI gate pass on a clean machine & Met: Compose stack run on a team laptop (16~October); both CI workflows pass on GitHub-hosted runners \\
\end{xltabular}}
